% GENERATED FILE. Edit canonical lessons, then run npm run generate:books.
% canonical-source-hash: fnv1a64:7e8de46630a25f98
% canonical-lessons: PT-C93-tapete, PT-C93-televisao, PT-C93-computador, PT-C93-frigorifico, PT-C93-forno

\chapter{Rug, Television, Computer, Fridge, Oven}
\label{ch:pt-rug-television-computer-fridge-oven}
\hlchaptermodality{93}

\begin{chapteropening}
\textbf{By the end of this chapter:} \emph{I can say a rug, a television, a computer, a fridge and an oven.}

Complete the last lesson of chapter 93: I can say a rug, a television, a computer, a fridge and an oven.
\end{chapteropening}

\section[o tapete]{o tapete — a rug, a carpet}

\label{lesson:PT-C93-tapete}



\begin{quote}
Before the new one: say the Portuguese for a bus, then the Portuguese for a train.
\end{quote}

\subsection*{You'll want to know: o tapete}
\textbf{o tapete} — \textquotedblleft{}a rug, a carpet\textquotedblright{}.

\begin{cousinweb}
\textbf{tapete} $\leftarrow$ Latin \emph{tapete}, \textquotedblleft{}carpet, rug\textquotedblright{}, from Greek \emph{tapēs} (Lewis \& Short). EtymDB 2.1 traces it only to Old Portuguese \emph{tapete}.
\end{cousinweb}

\begin{grammarlens}[title={What you've built}]
One more word for this chapter.
\end{grammarlens}

\subsection*{Your turn}
\begin{itemize}
\raggedright
  \item \emph{Say it:} \emph{o tapete}
  \item \emph{Say it:} \emph{o tapete}, once more
  \item \emph{Say it:} say \emph{o tapete}
  \item \emph{Recall:} say the Portuguese for a bus, then the Portuguese for a train, then say \emph{o tapete} again
\end{itemize}

\begin{tcolorbox}[breakable,colback=teal!4,colframe=teal!35!black,title={Before you move on}]
What does o tapete mean? (\textquotedblleft{}A rug, a carpet\textquotedblright{}.) Say it once more.
\end{tcolorbox}
\section[a televisão]{a televisão — a television}

\label{lesson:PT-C93-televisao}



\begin{quote}
Before the new one: say the Portuguese for a train, then the Portuguese for a rug.
\end{quote}

\subsection*{You'll want to know: a televisão}
\textbf{a televisão} — \textquotedblleft{}a television\textquotedblright{}.

\begin{cousinweb}
\textbf{televisão} is built from \emph{tele-}, \textquotedblleft{}far\textquotedblright{} (from Greek \emph{tēle}), and the word for vision (EtymDB 2.1), which $\leftarrow$ Latin \emph{visio}, \textquotedblleft{}sight\textquotedblright{}, from \emph{videre}, \textquotedblleft{}to see\textquotedblright{}.
\end{cousinweb}

\begin{grammarlens}[title={What you've built}]
One more word for this chapter.
\end{grammarlens}

\subsection*{Your turn}
\begin{itemize}
\raggedright
  \item \emph{Say it:} \emph{a televisão}
  \item \emph{Say it:} \emph{a televisão}, once more
  \item \emph{Say it:} say \emph{a televisão}
  \item \emph{Recall:} say the Portuguese for a train, then the Portuguese for a rug, then say \emph{a televisão} again
\end{itemize}

\begin{tcolorbox}[breakable,colback=teal!4,colframe=teal!35!black,title={Before you move on}]
What does a televisão mean? (\textquotedblleft{}A television\textquotedblright{}.) Say it once more.
\end{tcolorbox}
\section[o computador]{o computador — a computer}

\label{lesson:PT-C93-computador}



\begin{quote}
Before the new one: say the Portuguese for a rug, then the Portuguese for a television.
\end{quote}

\subsection*{You'll want to know: o computador}
\textbf{o computador} — \textquotedblleft{}a computer\textquotedblright{}.

\begin{cousinweb}
\textbf{computador} $\leftarrow$ Latin \emph{computator}, \textquotedblleft{}reckoner\textquotedblright{} (EtymDB 2.1), from \emph{computare}, \textquotedblleft{}to count up\textquotedblright{} — a learned borrowing from Latin.
\end{cousinweb}

\begin{grammarlens}[title={What you've built}]
One more word for this chapter.
\end{grammarlens}

\subsection*{Your turn}
\begin{itemize}
\raggedright
  \item \emph{Say it:} \emph{o computador}
  \item \emph{Say it:} \emph{o computador}, once more
  \item \emph{Say it:} say \emph{o computador}
  \item \emph{Recall:} say the Portuguese for a rug, then the Portuguese for a television, then say \emph{o computador} again
\end{itemize}

\begin{tcolorbox}[breakable,colback=teal!4,colframe=teal!35!black,title={Before you move on}]
What does o computador mean? (\textquotedblleft{}A computer\textquotedblright{}.) Say it once more.
\end{tcolorbox}
\section[o frigorífico]{o frigorífico — a fridge}

\label{lesson:PT-C93-frigorifico}



\begin{quote}
Before the new one: say the Portuguese for a television, then the Portuguese for a computer.
\end{quote}

\subsection*{You'll want to know: o frigorífico}
\textbf{o frigorífico} — \textquotedblleft{}a fridge\textquotedblright{}. In Brazil, \textbf{a geladeira}.

\begin{cousinweb}
\textbf{frigorífico} $\leftarrow$ Latin \emph{frigorificus}, \textquotedblleft{}cooling\textquotedblright{} (Lewis \& Short), from \emph{frigus}, \textquotedblleft{}cold\textquotedblright{}, and \emph{facere}, \textquotedblleft{}to make\textquotedblright{} — a learned borrowing from Latin.
\end{cousinweb}

\begin{grammarlens}[title={What you've built}]
One more word for this chapter.
\end{grammarlens}

\subsection*{Your turn}
\begin{itemize}
\raggedright
  \item \emph{Say it:} \emph{o frigorífico}
  \item \emph{Say it:} \emph{o frigorífico}, once more
  \item \emph{Say it:} say \emph{o frigorífico}
  \item \emph{Recall:} say the Portuguese for a television, then the Portuguese for a computer, then say \emph{o frigorífico} again
\end{itemize}

\begin{tcolorbox}[breakable,colback=teal!4,colframe=teal!35!black,title={Before you move on}]
What does o frigorífico mean? (\textquotedblleft{}A fridge\textquotedblright{}.) Say it once more.
\end{tcolorbox}
\section[o forno]{o forno — an oven}

\label{lesson:PT-C93-forno}



\begin{quote}
Before the new one: say the Portuguese for a computer, then the Portuguese for a fridge.
\end{quote}

\subsection*{You'll want to know: o forno}
\textbf{o forno} — \textquotedblleft{}an oven\textquotedblright{}.

\begin{cousinweb}
\textbf{forno} $\leftarrow$ Latin \emph{furnus}, \textquotedblleft{}oven\textquotedblright{} (EtymDB 2.1) — inherited.
\end{cousinweb}

\begin{grammarlens}[title={What you've built}]
One more word for this chapter.
\end{grammarlens}

\subsection*{Your turn}
\begin{itemize}
\raggedright
  \item \emph{Say it:} \emph{o forno}
  \item \emph{Say it:} \emph{o forno}, once more
  \item \emph{Say it:} say \emph{o forno}
  \item \emph{Recall:} say the Portuguese for a computer, then the Portuguese for a fridge, then say \emph{o forno} again
\end{itemize}

\begin{tcolorbox}[breakable,colback=teal!4,colframe=teal!35!black,title={Before you move on}]
What does o forno mean? (\textquotedblleft{}An oven\textquotedblright{}.) Say it once more.
\end{tcolorbox}
